\chapter*{\RL{ملخص}}
\phantomsection
\addcontentsline{toc}{chapter}{Arabic Abstract}

\setlength{\emergencystretch}{3em}
\begin{RLtext}

\noindent يقدّم هذا التقرير مشروع نهاية الدراسات المنجز بشركة إس تي مايكروإلكترونيكس بتونس.

ويقارن منظومة \LR{STM32Cube} بمنظومات متحكمات دقيقة منافسة. في مهمة مستقلة، زودنا فريق الإشراف بقائمة من لوحات المنافسين وطلب منا تحديد أقرب لوحة \LR{STM32} لكل منها. لم تحدد هذه الدراسة لوحات حملات القياس اللاحقة. بدأ العمل التجريبي بعد ذلك بدراسة \LR{USB HID} و\LR{CDC} و\LR{MSC} على منصات إس تي ورينيساس و\LR{NXP}. خفّض تحسين مشروع إس تي المرجعي استهلاك \LR{RAM} وزمن التهيئة في سيناريوهات رينيساس الثلاثة، بينما ارتبط تفسير نتائج \LR{NXP} بحدود القياس وبإدراج كود \LR{RCC} أو استبعاده.

اعتمدنا وظائف متكافئة وشروط بناء موثقة، واستخرجنا البصمة الذاكرية من ملفات الربط، كما راجعنا التكافؤ الدلالي وعدالة المقارنة.

طوّرنا أداتين مساعدتين: \LR{MCU Workflow Studio} لتنظيم تحليل ملفات الربط بين المصنّعين، و\LR{MCU Benchmark Copilot Agent} لتنظيم إنشاء المشاريع ونقلها وبنائها وتتبع مصادرها والتحقق منها. تساعد الأداتان على توثيق سير العمل، لكن صحة كل نتيجة تبقى مرتبطة بأدلة البناء والقياس والعتاد.

غطّت مقارنة \LR{ST--GD32} سيناريوهات \LR{USB HID} و\LR{CDC} وثلاثة سيناريوهات لـ\LR{FreeRTOS}. استهلك متحكم إس تي ذاكرة فلاش أكثر بنسبة 35--88\%، واستخدم ذاكرة \LR{RAM} أقل بنسبة 57--65\% في سيناريوهات \LR{USB}. أظهرت القياسات المسجلة فرقاً قدره 61,7\% في زمن تهيئة \LR{USB} لصالح إس تي، غير أن غياب العينات المتكررة وحدود القياس المفصلة يمنع تعميم تفسير سببي.

جمعت مقارنة \LR{ST--Texas Instruments} بين تحليل ثابت للمشغلات واثني عشر مشروعاً متناظراً. استخدم متحكم إس تي فلاشاً إجمالياً أكثر بنسبة 2,8\% وذاكرة \LR{RAM} أقل بنسبة 5,0\%. وفي سيناريوهات \LR{Free\-RTOS} الثلاثة، استخدمت صور إس تي فلاشاً أقل بنسبة 27,9--29,6\%، بينما وفّر متحكم \LR{MSPM33} ما بين 56 و76 بايتاً من \LR{RAM}. جُمعت أيضاً قياسات زمن التنفيذ، لكن إس تي عمل بتردد 144 ميغاهرتز مقابل 32 ميغاهرتز لـ\LR{MSPM33}. أصبح السيناريو الأساسي متقارباً بعد التطبيع حسب التردد، في حين بقي تفسير السيناريوهين التعاوني والاستباقي محدوداً باختلاف إعدادات النواة وآليات التزامن. لذلك لا تسمح الأدلة الحالية باستنتاج عام حول كفاءة المجدول أو استهلاك الطاقة.

\end{RLtext}

\noindent\rule[2pt]{\textwidth}{0.5pt}

\begin{RLtext}

{\textbf{الكلمات المفتاحية:}}

\LR{STM32Cube}، قياس الأداء، البصمة الذاكرية، متحكم دقيق، \LR{FreeRTOS}، \LR{USB}، \LR{GD32}، \LR{Texas Instruments}، مقارنة بين المصنّعين.
\\

\end{RLtext}

\noindent\rule[2pt]{\textwidth}{0.5pt}